\section{Why small rounding differences become different answers}
\label{sec:mechanism}

Section~\ref{sec:background} explains why execution-equivalent configurations produce different
float32 results: different kernels sum in different orders, and only some tiers fuse multiply
and add. That alone does not explain the size of the effect. Float32 reordering perturbs a
result by roughly $10^{-7}$ relative per operation, yet Table~\ref{tab:kl} shows mean KL
divergences of $10^{-3}$ to $10^{-2}$ and top-1 disagreement on several percent of positions.
This section tests one explanation and checks the code-level predictions.

\subsection{Activation quantization amplifies rounding differences}
Before every quantized matrix product, llama.cpp converts the float32 activation row into 8-bit
integers per block: blocks of 32 values for Q8\_0 (used with Q4\_0 and Q8\_0 weights) and of 256
values for Q8\_K (used with K-quants). Each value is mapped to
$\mathrm{round}(127\,x/\max|x_{\text{block}}|)$. A perturbation of one unit in the last place
almost never matters, but when it moves a value across a rounding boundary, the stored integer
changes by one, which is an error of $\max|x|/127$, several orders of magnitude larger than the
perturbation that caused it. A small simulation (\texttt{amp\_sim.py}) makes the scale concrete:
perturbing Gaussian activations with sparse outliers by one ulp of relative error changes only
$4.8\times10^{-6}$ of the Q8\_0 codes, but those changes raise the relative error of the
quantized row from $1.2\times10^{-7}$ to $2.9\times10^{-5}$, a factor of about 250, in a single
matmul. A transformer applies this step seven times per layer, and the errors feed forward.

Weights stored as F16 also quantize activations, but to F16 (11 significant bits instead of
about 8, and no shared block scale), and F32 weights with an F32 KV cache do not discretize
activations at all. If amplification is the main mechanism, the same execution change should
produce far smaller divergence for those formats. I pre-registered this test after one probe
(Addendum~A in \texttt{PREDICTIONS.md}) and ran it on Qwen3-0.6B in six formats
(Table~\ref{tab:amp}).

\begin{table}[t]
\centering\small
\caption{Divergence from the \texttt{alderlake} reference for the same model in six formats,
under two ISA changes and a thread-count change (8 chunks of 512 tokens). ``$\le10^{-6}$'' is
the storage floor of the measurement: configurations that are bitwise identical also read
$|\text{KLD}|\le10^{-6}$. F32+f32KV also stores the KV cache in F32.}
\label{tab:amp}
\begin{tabular}{lrrrrrr}
\toprule
format & \multicolumn{2}{c}{ivybridge} & \multicolumn{2}{c}{sse42} & \multicolumn{2}{c}{t4}\\
 & mean KLD & top-1 \% & mean KLD & top-1 \% & mean KLD & top-1 \%\\
\midrule
F32+f32KV & $\le 10^{-6}$ & 100.00 & $\le 10^{-6}$ & 100.00 & $\le 10^{-6}$ & 100.00\\
F32 & $\le 10^{-6}$ & 99.95 & $\le 10^{-6}$ & 99.95 & $\le 10^{-6}$ & 100.00\\
F16 & $\le 10^{-6}$ & 99.90 & $\le 10^{-6}$ & 99.90 & $\le 10^{-6}$ & 100.00\\
Q8\_0 & 0.0013 & 98.23 & 0.0014 & 98.04 & $\le 10^{-6}$ & 100.00\\
Q4\_K\_M & 0.0121 & 93.82 & 0.0118 & 93.68 & $\le 10^{-6}$ & 100.00\\
Q4\_0 & 0.0024 & 97.11 & 0.0026 & 97.16 & $\le 10^{-6}$ & 100.00\\
\bottomrule
\end{tabular}

\end{table}

The prediction held. With nothing discretized (F32 weights, F32 KV cache), removing FMA
(\texttt{ivybridge}) or SIMD (\texttt{sse42}) changes no top-1 token in 3,072 positions and the
KL divergence is at the measurement floor. Storing the KV cache in F16 or the weights in F16
changes one to three positions. Switching to Q8\_0 weights, with everything else equal, raises
the mean KLD from below $10^{-6}$ to $1.3\times10^{-3}$ and the top-1 disagreement to about 2\%.
Q4\_0 is about twice Q8\_0, and Q4\_K\_M, whose activations use 256-value Q8\_K blocks, is about
ten times Q8\_0. I did not isolate why Q4\_K\_M is so much more sensitive; the larger activation
block (a coarser step for most values in the block) is one candidate. The thread-count control
is exactly zero in every format.

The same amplification explains why the AVX-only tiers differ so much from the AVX2 tiers even
though their integer arithmetic is identical: any float32 difference upstream of an activation
quantizer, whether from FMA, a different summation order or a different kernel, is converted
into integer differences of a full step at the next layer.

\subsection{Code-level predictions}
Table~\ref{tab:pred} lists the predictions written before the runs and their outcomes.

\begin{table}[t]
\centering\small
\caption{Pre-registered predictions (\texttt{PREDICTIONS.md}) and outcomes.}
\label{tab:pred}
\begin{tabular}{p{0.37\linewidth}p{0.13\linewidth}p{0.40\linewidth}}
\toprule
prediction & outcome & evidence\\
\midrule
P1 threads 4/8/12 identical & supported & 6/6 models (Table~\ref{tab:kl})\\
P2 \texttt{haswell} = \texttt{alderlake} & supported & 6/6 models\\
P3 \texttt{ivybridge} differs & supported & 6/6 models\\
P4 \texttt{sandybridge} differs from \texttt{ivybridge}; \texttt{sse42} differs from AVX tiers & not tested & both differ from $R$; the tiers were never compared with each other\\
P5 micro-batch 1 differs & supported & 6/6 models\\
P5 micro-batch 7 differs & \textbf{falsified} for 5/6 & differs only for Q4\_K\_M; identical for Q4\_0 (3 models), Q8\_0, F16\\
P6 repack off differs for Q4\_0, Q4\_K\_M only & supported & differs for all four Q4 files, identical for Q8\_0 and F16\\
P7 flash attention differs & supported & 6/6 models\\
P8 rerun of $R$ identical & supported & 980/980 items bitwise equal\\
P9 prompt cache changes some items & supported & 68\% of items not bitwise equal (Q4\_K\_M)\\
P10 two identical parallel-server runs differ & supported & Q4\_K\_M only (the only model run twice)\\
P11 flips concentrate on low-margin items & supported & Figure~\ref{fig:margins}\\
A1 F32+f32KV $<$ F16 $<$ Q8\_0, $\ge$10$\times$ & supported & Table~\ref{tab:amp}\\
A2 thread control exactly zero & supported & all six formats\\
\bottomrule
\end{tabular}

\end{table}

Of the thirteen predictions the runs tested, twelve held. The clear failure is micro-batch 7, which I
expected to differ from micro-batch 512 for every model because the repacked kernels process
the rows of a micro-batch in groups of four with one routine and the remainder with another
(F4 in \texttt{PREDICTIONS.md}). It differs only for Q4\_K\_M. For Q4\_0 the grouped and
remainder routines are generated from one template and evidently produce the same bits per row,
while the Q4\_K routines are separately written kernels that do not. For Q8\_0, which is not
repacked, micro-batch 7 also moves the matrix products from the tinyBLAS path to the per-row
dot product path (F5), yet the logits are identical; both paths accumulate each 32-value block
with one fused multiply-add into an 8-lane float32 accumulator and reduce it at the end, which is
consistent with the match, though I did not check them instruction by instruction. Micro-batch
1 differs everywhere, so the attention path, which changes with the number of query rows, is the
likely source of that difference. The general lesson is that whether a performance knob changes
the numbers depends on the weight format, so a setting found safe for one file is not safe for
another.
